\documentclass[10pt,a4paper]{amsart}
\usepackage[margin=19mm]{geometry}
\usepackage{amsmath,amssymb,mathtools}
\usepackage[scr=rsfs,cal=euler]{mathalfa}
\usepackage{xfrac}
\usepackage[unicode,hidelinks,bookmarks=false]{hyperref}
\newcommand{\ie}{i.e.\ }
\newcommand{\cf}{cf.\ }
\newcommand{\resp}{respectively\ }
\newcommand{\Subsection}{\subsection}
\providecommand{\nobreakdash}{\nobreak\hbox{-}}
\setlength{\emergencystretch}{2em}
\multlinegap0pt
\usepackage{enumerate}
\usepackage{amssymb}
\usepackage{dsfont}
\usepackage{mathtools}
\usepackage{xcolor}
\usepackage{tikz}
\newtheorem{prop}{Proposition}
\newtheorem{lem}{Lemma}
\newtheorem{definition}{Definition}
\newcommand{\F}{\mathbb{F}}
\title{Moments of traces of random symplectic matrices and hyperelliptic $L$-functions}
\author{Alexei Entin, Noam Pirani}
\date{Source: arXiv:2409.04844v2 (CC BY 4.0)}
\begin{document}
\maketitle

\noindent\emph{Source and licence.} Moments of traces of random symplectic matrices and hyperelliptic $L$-functions. Version arXiv:2409.04844v2 (CC BY 4.0), distributed by arXiv under the Creative Commons Attribution 4.0 International licence, http://creativecommons.org/licenses/by/4.0/, as recorded in the arXiv licence metadata for this exact version and shown on the abstract page https://arxiv.org/abs/2409.04844v2. Full licence notice: \texttt{licenses/arxiv-2409.04844-CC-BY-4.0.txt}.\par\smallskip

\noindent\textit{Editorial scope.} Counted unit: the article's prop prop\_main\_sum\_estimate (source0.tex lines 554--560). The article's complete original proof of it is delivered with it (source0.tex lines 562--635, one \texttt{proof} environment of 74 lines). No mathematics is written, repaired or re-derived: the statement and the proof are the article's own lines, in the article's own notation, and no retained line is substituted or re-flowed.  Retained uncounted as the source-local context the proof needs: lines 166--167; lines 278--280; lines 283--290; lines 345--347; lines 485--491; lines 526--526; lines 527--531; lines 533--535.  Nothing was omitted from any retained span, so the package carries no omission record.  No retained line contains a citation, so the wrapper carries no bibliography and no bibliography entry.  No retained line contains a figure environment, an image-inclusion command or drawing code, so no image is delivered and the package declares no asset dependency.  Editorial formatting only, in addition: an AMS \texttt{amsart} preamble that restates the article's own declarations and macros used by the retained lines, namely the environments \texttt{prop}, \texttt{lem}, \texttt{definition} on separate counters, together with the standard packages those lines need.  The retained environments print in this document as Proposition 1, Lemma 1, Definition 1 in order of appearance, whereas the article prints the same items with the numbering of its own full document.  The complete native source of the article as distributed in the arXiv e-print for this exact version is bundled unchanged as \texttt{sources/164927/source0.tex}.
\begin{equation}\label{eq: def phi}\phi(n,\mathbf c)=\left[\begin{array}{ll}
-\sum_{\mathbf d\le\mathbf c\atop{|\mathbf d|\le{|\mathbf c|}/2-n-1}}(-1)^{\ell(\mathbf d)}{\mathbf c\choose\mathbf d},&|\mathbf c|>0,\,2\mid |\mathbf c|,\\1,&|\mathbf c|=0, \\0,&|\mathbf c|>0,\,2\nmid|\mathbf c|. \end{array}\right.\end{equation}

\begin{equation}\label{eq: def T}
T(n;\mathbf a)=\sum_{h\,\mathrm{monic}\atop{\deg h=2n+1}}\sum_{(P_{ji}:\,1\le j\le k,\,1\le i\le a_j)\atop{\text {distinct primes}\atop{\deg P_{ji}=j}}}\prod_{1\le j\le k\atop{1\le i\le a_j}}\Lambda(P_{ji})\left(\frac h{P_{ji}}\right).
\end{equation}

\begin{prop}\label{prop_S_sum_estimate}
Assume $|\mathbf a|\le 4n+1$. Then
\begin{equation*}
    q^{-2n-1-|\mathbf a|/2}T(n;\mathbf a)=\phi(n,\mathbf a)+O\left(q^{-1/2}\right),
\end{equation*}
where $\phi$ is defined by (\ref{eq: def phi}).

\end{prop}

\begin{lem}\label{lem: weil gen} Let $h\in\F_q[x]$ be a monic polynomial of degree $2n+1$, $r,j_1,\ldots,j_r$ fixed natural numbers. Then
$$\sum_{P_1,\ldots,P_r\atop{\mathrm{distinct}\atop{\deg P_l=j_l}}}\left(\frac h{P_1\cdots P_r}\right)=O\left(q^{\frac{j_1+\cdots+j_r}2}\right).$$
\end{lem}

\begin{prop}\label{prop_sum_to_pair_up}
Let $\mathbf b$ be a partition supported on $[1,k]$. Then,

\begin{equation}\label{eq: contribution of squares} q^{-|\mathbf b|/2}\sum_{(Q_{ji}:\,1\le j\le k,\,1\le i\le b_j)\atop{\prod Q_{ji}=\square}}\prod_{1\le j\le k\atop{1\le i\le b_j}}\Lambda(Q_{ji})=g(\mathbf b)+O(q^{-1}),
\end{equation}
where $\prod Q_{ji}$ denotes the product over $1\le j\le k,1\le i\le b_j$.
\end{prop}

Denote \begin{equation}\label{eq: def A}A=\{ji:\,1\le j\le k,\,1\le i\le a_j\}.\end{equation}

\begin{definition}\label{def: compatible}Let $K\subset A$ be a subset and $\sim$ an equivalence relation on $A\setminus K$. We say that a collection $(Q_{ji}:\,1\le j\le k,\,1\le i\le a_j)$ is \emph{compatible with} $\sim$ if the following conditions hold:
\begin{enumerate}\item[(i)] $Q_{ji}=\text{prime}^2$ iff $ji\in K$ (so $Q_{ji}$ is prime iff $ji\in A\setminus K$).
\item[(ii)] If $ji,j'i'\in A\setminus K$ then $Q_{ji}=Q_{j'i'}$ iff $ji\sim j'i'$ (in particular $j=j'$ if $ji\sim j'i'$).
\end{enumerate}
\end{definition}

\begin{lem}\label{lem_sums_in_inclusion_exclusion} Let $K\subset A$ be a subset and $\sim$ an equivalence relation on $A\setminus K$ such that at least one equivalence class has size $\ge 3$. Then
\begin{equation}\label{eq: lem triple}q^{-2n-1-|\mathbf a|/2}\sum_{\deg h=2n+1\atop{\text{monic}}}\sum_{(Q_{ji}:\,1\le j\le k,\,1\le i\le a_j)\atop{\text{compatible with $\sim$}}}\prod_{1\le j\le k\atop{1\le i\le a_j}}\Lambda(Q_{ji})\left(\frac h{Q_{ji}}\right)=O(q^{-1}).\end{equation}
\end{lem}

\begin{prop}\label{prop_main_sum_estimate}
    Let $\mathbf a$ be supported on $[1,k]$, and assume that $|\mathbf a|\le 4n+1$. Then
    \begin{equation*}
    R:=q^{-2n-1-|\mathbf a|/2}\sum_{\deg h=2n+1\atop{\mathrm{monic}}}\sum_{(Q_{ji}:\,1\le j\le k,\,1\le i\le a_j)}\prod_{1\le j\le k\atop{1\le i\le a_j}}\Lambda(Q_{ji})\left(\frac{h}{Q_{ji}}\right)
    =\sum_{\mathbf b\le\mathbf a}{\mathbf a\choose\mathbf b}g(\mathbf b)\phi(n,\mathbf a-\mathbf b)+O(q^{-{1/2}}).
    \end{equation*}
\end{prop}

\begin{proof} In what follows $\subset$ denotes non-strict inclusion of sets. We split the sum on the LHS of the proposition into terms in the following way: for every choice of subset $J\subset A$ of the indices (recall that $A$ is defined by (\ref{eq: def A})), we say that a choice $(Q_{ji}:1\le j\le k,1\le i\le a_j)$ respects $J$ if
\begin{enumerate}
    \item[(i)] $\prod_{ji\in J}Q_{ji}=\square$,
    \item[(ii)] $Q_{ji}$ for $ji\notin J$ are distinct primes.
\end{enumerate} We now consider the sums 
$$
R_J=q^{-2n-1-|\mathbf a|/2}
\sum_{\deg h=2n+1\atop{\mathrm{monic}}}
\sum_
{(Q_{ji}:\,1\le j\le k,\,1\le i\le a_j)\atop{\text{respects } J}
}
\prod_{1\le j\le k\atop{1\le i\le a_j}}\Lambda(Q_{ji})\left(\frac{h}{Q_{ji}}\right).
$$
We also define for subsets $J_1,\ldots,J_s\subset A$,
$$
R_{J_1,\ldots,J_s}=q^{-2n-1-|\mathbf a|/2}
\sum_{\deg h=2n+1\atop{\mathrm{monic}}}
\sum_
{(Q_{ji}:\,1\le j\le k,\,1\le i\le a_j)\atop{\text{respects } J_1,\ldots,J_s}
}
\prod_{1\le j\le k\atop{1\le i\le a_j}}\Lambda(Q_{ji})\left(\frac{h}{Q_{ji}}\right).
$$
Since any choice of $Q_{ji}$ respects some subset $J$ (collect all the squares and pairs of repeating primes so that only distinct primes remain), by inclusion-exclusion we have 
\begin{equation}\label{eq: inclusion-exclusion}
R =\sum_{J_1}R_{J_1}-\sum_{J_1,J_2}R_{J_1,J_2}+\sum_{J_1,J_2,J_3}R_{J_1,J_2,J_3}-\ldots,
\end{equation}
where summation in each term is over unordered tuples of distinct subsets of $A$.
We will show that $R_{J_1,\ldots,J_s}=O(q^{-1})$ for $s\ge 2$ (and $J_1,\ldots,J_s$ distinct) and evaluate $R_J$ for every $J$, which will complete the evaluation of $R$ (up to an $O(q^{-1})$ error term).


Let $J_1,\ldots,J_s,\,s\ge 2$ be distinct. Consider a collection $(Q_{ji}:\,1\le j\le k,\,1\le i\le a_j)$. Denote $K=\{ji: Q_{ji}=\text{prime}^2\}$ and define an equivalence relation $\sim$ on $A\setminus K$ (which consists of the prime $Q_{ji}$) by $ji\sim j'i'$ iff $Q_{ji}=Q_{j'i'}$. This is the unique choice of $(K,\sim)$ that is compatible with $(Q_{ji})$ in the sense of Definition \ref{def: compatible}. The collection $(Q_{ji})$ respects $J_l$ iff the following conditions hold:
\begin{enumerate}
\item[(i)] $K\subset J_l$.
\item[(ii)] For each equivalence class $C$ of $\sim$ we have that $|C\cap J_l|$ is even and $|C\setminus J_l|\le 1$.
\end{enumerate}
Therefore the condition that $(Q_{ji})$ respects $J_1,\ldots,J_s$ depends only on $(K,\sim)$. Moreover if (i-ii) above hold with two distinct $J_1,J_2$ then one of the equivalence classes $C$ must be of size $\ge 3$. To see this let (WLOG) $uv\in J_1\setminus J_2$. Then $Q_{uv}$ is a prime which appears at least twice in $(Q_{ji})_{ji\in J_1}$ (because $\prod_{ji\in J_1}Q_{ji}=\square$) and therefore at least once in $(Q_{ji})_{ji\in J_2}$ (because $(Q_{ji})_{ji\in A\setminus J_2}$ are distinct primes). It follows that $Q_{uv}$ appears at least twice in $(Q_{ji})_{ji\in J_2}$ (because $\prod_{ji\in J_2}Q_{ji}=\square$) and since $uv\not\in J_2$ it appears $\ge 3$ times in $(Q_{ji})_{ji\in A}$. Its equivalence class $C$ has therefore size $\ge 3$. 

Using Lemma \ref{lem_sums_in_inclusion_exclusion} we obtain
$$R_{J_1,\ldots,J_s}=\sum_{K,\sim}q^{-2n-1-|\mathbf a|/2}\sum_{\deg h=2n+1\atop{\text{monic}}}\sum_{(Q_{ji}:\,1\le j\le k,\,1\le i\le a_j)\atop{\text{compatible with $\sim$}}}\prod_{1\le j\le k\atop{1\le i\le a_j}}\Lambda(Q_{ji})\left(\frac h{Q_{ji}}\right)=O(q^{-1}).$$
Here the sum $\sum_{K,\sim}$ is over all $(K,\sim)$ such that compatibility with it implies respecting $J_1,\ldots,J_s$.




We conclude that $R_{J_1,\ldots,J_s}=O(q^{-1})$ and from (\ref{eq: inclusion-exclusion}) that
\begin{equation}\label{eq: R through RJ}R=\sum_{J\subset A}R_J+O(q^{-1}).\end{equation}
Next we evaluate $R_J$ for $J\subset A$.
With each $J\subset A$ we associate a partition $\mathbf b\le\mathbf a$ by $b_j=|\{1\le i\le a_j: ji\in J\}|$.
The number of subsets $J$ corresponding to a given $\mathbf b\le\mathbf a$ is exactly $\mathbf a\choose\mathbf b$ and $R_J$ depends only on $\mathbf b$ and not on the specific $J$. Hence from (\ref{eq: R through RJ}) we have
\begin{equation}\label{eq: R through Rb}R=\sum_{\mathbf b\le\mathbf a}{\mathbf a\choose\mathbf b}R_{\mathbf b}+O(q^{-1}),\end{equation}
where
$$R_{\mathbf b}=q^{-2n-1-|\mathbf a|/2}
\sum_{\deg h=2n+1\atop{\mathrm{monic}}}
\sum_
{(Q_{ji}:\,1\le j\le k,\,1\le i\le b_j)\atop{\prod_{1\le j\le k\atop{1\le i\le b_j}}Q_{ji}=\square\atop{(Q_{ji})_{1\le j\le k\atop{b_j<i\le a_j}}\text{distinct primes}}}}
\prod_{1\le j\le k\atop{1\le i\le a_j}}\Lambda(Q_{ji})\left(\frac{h}{Q_{ji}}\right)
$$
(the latter expression is just $R_J$ for $J=\{ji:\,1\le j\le k,\,1\le i\le b_j\}$, which corresponds to $\mathbf b$).
It remains to evaluate $R_{\mathbf b}$ for each $\mathbf b\le\mathbf a$.

Fix a partition $\mathbf b\le \mathbf a$. Let $(Q_{ji}:\,1\le j\le k,\,1\le i\le a_j)$ be a choice respecting $J=\{ji:\,1\le j\le k,\,1\le i\le b_j\}.$ Recall that this means that $\prod_{1\le j\le k\atop{1\le i\le b_j}}Q_{ji}=\square$ and $Q_{ji},\,1\le j\le k,\,b_j< i\le a_j$ are distinct primes. If we fix $Q_{ji}$ for $1\le j\le k,\,1\le i\le b_j$ then 
\begin{multline*}
\sum_{\deg h=2n+1\atop{\text{monic}}}\sum_{(Q_{ji}:\,1\le j\le k,\,b_j<i\le a_j)\atop{\text{distinct primes}}}\prod_{1\le j\le k\atop{1\le i\le a_j}}\Lambda(Q_{ji})\left(\frac{h}{Q_{ji}}\right)=
\\=\prod_{1\le j\le k\atop{1\le i\le b_j}}\Lambda(Q_{ji})\sum_{\deg h=2n+1\atop{\text{monic}\atop{\left(h,\prod_{1\le j\le k,\,1\le i\le b_j} Q_{ji}\right)=1}}}\sum_{(Q_{ji}:\,1\le j\le k,\,b_j<i\le a_j)\atop{\text{distinct primes}}}
\prod_{1\le j\le k\atop{b_j<i\le a_j}}\Lambda(Q_{ji})\left(\frac{h}{Q_{ji}}\right)=\\
=\prod_{1\le j\le k\atop{1\le i\le b_j}}\Lambda(Q_{ji})\sum_{\deg h=2n+1\atop{\text{monic}}}\sum_{(Q_{ji}:\,1\le j\le k,\, b_j<i\le a_j)\atop{\text{distinct primes}}}\prod_{1\le j\le k\atop{b_j<i\le a_j}}\Lambda(Q_{ji})\left(\frac{h}{Q_{ji}}\right)+O(q^{2n+|\mathbf a-\mathbf b|/2}).
\end{multline*}
When $\mathbf b$ is not the empty partition, the second equality holds because there are $O(q^{2n})$ polynomials $h$ which are not coprime with $\prod_{1\le j\le k\atop{1\le i\le b_j}}Q_{ji}$ and by Lemma \ref{lem: weil gen} each of them contributes $O(q^{|\mathbf a-\mathbf b|/2})$ to the sum. When $\mathbf b$ is the empty partition, the transition is immediate. Summing over the $O(q^{|\mathbf b|/2})$ appropriate choices of $(Q_{ji})_{1\le j\le k,\,1\le i\le b_j}$ (i.e. such that the product is a square) and using the notation (\ref{eq: def T}), Proposition \ref{prop_S_sum_estimate} and Proposition \ref{prop_sum_to_pair_up} (the propositions applied to $\mathbf a-\mathbf b$ and $\mathbf b$ respectively) we obtain
\begin{align*}
R_{\mathbf b}&=q^{-2n-1-|\mathbf a|/2}\left(\sum_{(Q_{ji}:\,1\le j\le k,\,1\le i\le b_j)\atop\prod Q_{ji}=\square}\prod_{1\le j\le k\atop{1\le i\le b_j}}\Lambda(Q_{ji})\right)\left(\sum_{\deg h=2n+1\atop{\text{monic}}}\sum_{(Q_{ji}:\,1\le j\le k,\,b_j<i\le a_j)\atop{\text{distinct primes}}}\prod_{1\le j\le k\atop{b_j<i\le a_j}}\Lambda(Q_{ji})\left(\frac{h}{Q_{ji}}\right)\right)\\&+O(q^{-1})
\\&=
q^{-2n-1-|\mathbf a|/2}\left(\sum_{(Q_{ji}:\,1\le j\le k,\,1\le i\le b_j)\atop\prod {Q_{ji}}=\square}\prod_{1\le j\le k\atop{1\le i\le b_j}}\Lambda(Q_{ji})\right)T(\mathbf a-\mathbf b)+O(q^{-1})
\\&=g(\mathbf b)\phi(n,\mathbf a-\mathbf b)+O(q^{-1/2}).\end{align*}
Combining this with (\ref{eq: R through Rb}) gives the assertion of Proposition \ref{prop_main_sum_estimate}.\end{proof}

\end{document}
